\section{Discussion}\label{sec:discussion}

\paragraph{Design implications.} For pointing interfaces on glasses, our results suggest five practical rules.
\begin{enumerate}
\item \emph{Choose the model by what the interface needs.} For a tight target box, a compact detector trained on pointing data is more accurate at the same compute. For naming the target, a structured model does as well, and it also provides the hand, the fingertip and the direction for feedback or gating in the same pass (Section~\ref{sec:choice}).
\item \emph{Train with distractor hands, whatever the model.} Pasting unlabelled hands into the training images makes SpatialNet and a generic detector equally stable when hands are pasted into the view, at no cost on clean images (Section~\ref{sec:design}).
\item \emph{Spend extra compute on resolution first.} A larger input gained more than a larger backbone (Section~\ref{sec:design}).
\item \emph{Confirm before acting.} Errors concentrate where similar objects lie side by side or the target is far from the fingertip, and single frames produce false alarms. Naming or highlighting the target before acting, and requiring several consecutive detections, protect the wearer from both (Section~\ref{sec:failures}).
\item \emph{Keep frames on the device.} A later recogniser needs only the target crop (Section~\ref{sec:choice}), which limits computation, bandwidth and the bystander imagery that leaves the glasses. Whether local inference also saves energy needs repeated measurements; the single observations of Section~\ref{sec:ondevice} only point in that direction.
\end{enumerate}

\paragraph{Limitations.} \emph{Data.} All accuracy figures except Section~\ref{sec:realframes} come from one dataset whose images are mostly landscape, partly composed or edited, at least 41\% derived from COCO, and in 4.7\% of test cases duplicated in the training or validation split (Appendix~\ref{app:data}); its demographic coverage is unknown. Frames are squashed to a square, whereas the glasses' camera is portrait. EgoPoint has no non-pointing hands and no pointing without a target, so the pointing and target gates are untested, and in deployment a non-pointing hand may be read as pointing. The model has no temporal context. \emph{Other hands.} Our distractors are synthetic and encode the prior that the wearer's arm enters from the bottom of the view; real bystander hands are untested. \emph{Evidence.} The ablations, the resolution runs and three of the four controlled detector runs are single training runs, and the per-component reading of the ablations is exploratory. We copied the large-model results rather than re-running them. The recognition test is category-level and uses zero-shot recognisers. The device results are single observations on a pre-production unit, and the detectors were not measured on the glasses. \emph{Interaction.} We measured the model, not the interaction: how wearers respond to its errors, and which confirmation designs suit them, needs a study with participants wearing the glasses.
